% ============================================================
% Round 09: GSPO
% ============================================================

\subsection{Round 09：GSPO}

\subsubsection{本轮问题}

\begin{conceptbox}
\textbf{学习目标：} 理解 GSPO 为什么要把 GRPO 的 token-level importance ratio 改成 sequence-level importance ratio，以及这如何提升 long-CoT / MoE RL 训练稳定性。\\
\textbf{主线位置：}
\[
\text{GRPO}
\to
\text{DAPO / Dr.GRPO}
\to
\text{GSPO}
\to
\text{Hybrid / Subsequence-level methods}
\]
\end{conceptbox}

\begin{conceptbox}
\textbf{Highlight：GSPO}\\
\textbf{论文：} \emph{Group Sequence Policy Optimization}\\
\textbf{arXiv：} \href{https://arxiv.org/abs/2507.18071}{2507.18071}，submitted 2025-07-24，last revised 2025-07-28\\
\textbf{Blog：} \href{https://qwenlm.github.io/blog/gspo/}{Qwen: GSPO: Towards Scalable Reinforcement Learning for Language Models}\\
\textbf{Code：} 论文未提供单独官方 repo；工程实现可参考 \href{https://github.com/modelscope/ms-swift}{modelscope/ms-swift} 和 \href{https://github.com/volcengine/verl}{volcengine/verl}\\
\textbf{关键词：} sequence-level importance ratio，sequence-level clipping，MoE RL stability，Routing Replay，Qwen3\\
\textbf{核心定位：} 把 GRPO 的 token-level policy optimization 改成 sequence-level policy optimization，让 reward 粒度和优化粒度对齐。
\end{conceptbox}

\subsubsection{GSPO 想修什么}

GRPO 的基础公式里，对每个 token 都有 importance ratio：

\[
\rho_{i,t}(\theta)
=
\frac{
\pi_\theta(y_{i,t}\mid x,y_{i,<t})
}{
\pi_{\theta_\text{old}}(y_{i,t}\mid x,y_{i,<t})
}
\]

然后用这个 token-level ratio 去更新每个 token：

\[
\min
\left(
\rho_{i,t}\hat{A}_i,\;
\operatorname{clip}(\rho_{i,t},1-\epsilon,1+\epsilon)\hat{A}_i
\right)
\]

但在 RLVR / reasoning RL 中，reward 通常不是 token-level 的，而是 sequence-level 的：

\[
r_i = R(x,y_i)
\]

例如数学题最后答案对了，整段回答拿 $1$；错了，整段回答拿 $0$。这时每个 token 都共享同一个回答级 advantage：

\[
\hat{A}_i
\]

这就产生了粒度不匹配：

\[
\text{reward 是整段级别}
\qquad
\text{ratio 却是逐 token 级别}
\]

\begin{summarybox}
\textbf{一句话：} GSPO 认为 GRPO 的 token-level ratio 会给 sequence-level reward 引入高方差噪声，因此应该把 ratio、clip、optimization 都提升到 sequence level。
\end{summarybox}

\subsubsection{为什么 Token-Level Ratio 会不稳}

在 long-CoT 中，一个回答可能有几千甚至上万 token。只要某些 token 的 ratio 特别大或特别小，就可能影响局部梯度和 clipping 行为。

例如同一个回答里，几个 token ratio 是：

\[
[1.02,\;0.98,\;5.00,\;1.01,\;0.20,\;1.03]
\]

这些极端 token ratio 不一定代表整段回答真的应该被剧烈更新，可能只是局部概率波动、数值误差、MoE routing 差异或训练/推理 engine 差异造成的。

但 GRPO 会在 token level 逐个处理：

\[
\rho_{i,t}
\to
\operatorname{clip}(\rho_{i,t})
\]

结果是：

\begin{itemize}[leftmargin=1.5em]
  \item 某些 token 被过度 clipping；
  \item 梯度噪声随序列长度累积；
  \item long-CoT 训练更容易不稳定；
  \item MoE 模型里 expert routing 波动会进一步放大问题。
\end{itemize}

\begin{intubox}
\textbf{直觉：} GRPO 像是对一篇作文里的每个字单独判断“新模型相对旧模型变了多少”；GSPO 认为 reward 明明是给整篇作文的，就应该从整篇作文的概率变化来判断。
\end{intubox}

\subsubsection{GSPO 的 Sequence-Level Ratio}

GSPO 对每个回答 $y_i$ 定义一个 sequence-level importance ratio：

\[
s_i(\theta)
=
\left[
\frac{
\pi_\theta(y_i\mid x)
}{
\pi_{\theta_\text{old}}(y_i\mid x)
}
\right]^{1/|y_i|}
\]

因为：

\[
\pi_\theta(y_i\mid x)
=
\prod_{t=1}^{|y_i|}
\pi_\theta(y_{i,t}\mid x,y_{i,<t})
\]

所以可以写成 log 形式：

\[
s_i(\theta)
=
\exp
\left(
\frac{1}{|y_i|}
\sum_{t=1}^{|y_i|}
\log
\frac{
\pi_\theta(y_{i,t}\mid x,y_{i,<t})
}{
\pi_{\theta_\text{old}}(y_{i,t}\mid x,y_{i,<t})
}
\right)
\]

注意这里有一个 $1/|y_i|$。它做的是长度归一化，避免整段概率连乘后随着长度指数级变小或 ratio 尺度随长度剧烈变化。

\begin{summarybox}
\textbf{GSPO ratio 的本质：}
\[
s_i(\theta)
=
\exp(\text{平均 token log-ratio})
\]
它不是看某一个 token 的概率变化，而是看整段回答平均意义上的概率变化。
\end{summarybox}

\subsubsection{GSPO 的目标函数}

GSPO 的 objective 是 sequence-level 的：

\[
\mathcal{J}_\text{GSPO}(\theta)
=
\mathbb{E}
\left[
\frac{1}{G}
\sum_{i=1}^{G}
\min
\left(
s_i(\theta)\hat{A}_i,\;
\operatorname{clip}(s_i(\theta),1-\epsilon,1+\epsilon)\hat{A}_i
\right)
\right]
\]

对比 GRPO：

\[
\mathcal{J}_\text{GRPO}(\theta)
=
\mathbb{E}
\left[
\frac{1}{G}
\sum_i
\frac{1}{|y_i|}
\sum_t
\min
\left(
\rho_{i,t}\hat{A}_i,\;
\operatorname{clip}(\rho_{i,t},1-\epsilon,1+\epsilon)\hat{A}_i
\right)
\right]
\]

核心差别：

\[
\text{GRPO：每个 token 一个 ratio，一个 clip}
\]

\[
\text{GSPO：每个 sequence 一个 ratio，一个 clip}
\]

\begin{warnbox}
\textbf{关键点：} GSPO 不是改掉 group advantage。它仍然可以使用同题 group relative advantage。它主要改的是 importance ratio 和 clipping 的粒度。
\end{warnbox}

\subsubsection{一个小例子}

假设某个回答有 $4$ 个 token，token-level ratio 是：

\[
[1.1,\;1.2,\;0.9,\;1.0]
\]

GRPO 会分别处理四个 ratio：

\[
\rho_{1},\rho_{2},\rho_{3},\rho_{4}
\]

GSPO 则先取平均 log-ratio：

\[
\frac{1}{4}
\left(
\log 1.1
+
\log 1.2
+
\log 0.9
+
\log 1.0
\right)
\]

再 exponentiate：

\[
s
=
\exp
\left[
\frac{1}{4}
(\log 1.1+\log 1.2+\log 0.9+\log 1.0)
\right]
\]

这等价于这些 ratio 的几何平均：

\[
s
=
(1.1 \times 1.2 \times 0.9 \times 1.0)^{1/4}
\]

\begin{intubox}
\textbf{直觉：} GSPO 用一整段回答的平均概率变化来决定更新强度，局部 token 的极端波动会被整体平均掉。
\end{intubox}

\subsubsection{为什么 GSPO 对 MoE 更重要}

MoE 模型每个 token 会选择不同 expert。GRPO 在 token level 比较：

\[
\pi_\theta(y_{i,t}\mid x,y_{i,<t})
\quad
\text{vs}
\quad
\pi_{\theta_\text{old}}(y_{i,t}\mid x,y_{i,<t})
\]

如果新旧模型的 expert routing 不一致，某些 token 的概率可能出现较大波动，导致 ratio 噪声变大。

为了解决这个问题，一些系统会用 Routing Replay：缓存 old policy 的 expert activation，在当前 policy 计算 ratio 时重放这些 routing patterns。

但 Routing Replay 会带来：

\begin{itemize}[leftmargin=1.5em]
  \item 额外显存开销；
  \item 额外通信和工程复杂度；
  \item 限制 MoE 模型使用完整容量。
\end{itemize}

GSPO 因为只关心 sequence-level likelihood，对单个 token 的 routing 波动不那么敏感，因此可以减少对 Routing Replay 这类工程 workaround 的依赖。

\begin{summarybox}
\textbf{MoE 场景下的直觉：} GRPO 的 token-level ratio 容易被 expert routing 局部波动干扰；GSPO 把比较提升到 sequence level，因此训练更稳定，基础设施也更简单。
\end{summarybox}

\subsubsection{GSPO 和 GRPO / DAPO / Dr.GRPO 的关系}

\begin{center}
{\small
\begin{tabular}{p{2.7cm}p{3.2cm}p{3.3cm}p{3.3cm}}
\hline
 & \textbf{GRPO} & \textbf{Dr.GRPO} & \textbf{GSPO} \\
\hline
核心问题
&
去 critic
&
去归一化偏置
&
去 token-level ratio 噪声
\\
Advantage
&
group relative
&
group mean baseline，去 std scaling
&
可继续用 group relative
\\
Ratio 粒度
&
token-level
&
通常仍是 token-level
&
sequence-level
\\
Clip 粒度
&
token-level clip
&
通常仍是 token-level clip
&
sequence-level clip
\\
主要收益
&
不用 critic
&
减少 length / difficulty bias
&
更稳定，尤其适合 long-CoT 和 MoE
\\
\hline
\end{tabular}
}
\end{center}

DAPO 和 GSPO 也可以对比：

\begin{itemize}[leftmargin=1.5em]
  \item DAPO 更像一套工程 recipe：dynamic sampling、Clip-Higher、token-level loss、overlong shaping；
  \item GSPO 更像一次优化粒度调整：把 token-level objective 改成 sequence-level objective。
\end{itemize}

\subsubsection{GSPO 的代价}

GSPO 不是没有 tradeoff。

它的优势是 sequence-level 信号更稳，但代价是 token-level credit assignment 变粗：

\[
\text{所有 token 共享同一个 sequence ratio}
\]

这可能导致：

\begin{itemize}[leftmargin=1.5em]
  \item 局部 token 的细粒度修正能力变弱；
  \item 对需要精细 token-level control 的任务不一定最优；
  \item 后续会出现 hybrid 方法，把 token-level 和 sequence-level ratio 结合起来。
\end{itemize}

\begin{warnbox}
\textbf{核心 tradeoff：} GRPO 更细，但噪声大；GSPO 更稳，但更粗。后续 hybrid / subsequence-level 方法就是在寻找中间粒度。
\end{warnbox}

\subsubsection{本轮最小结论}

\begin{summarybox}
\textbf{GSPO 的核心：}\\
把 GRPO 的 token-level ratio：
\[
\rho_{i,t}
=
\frac{\pi_\theta(y_{i,t}\mid x,y_{i,<t})}
{\pi_{\theta_\text{old}}(y_{i,t}\mid x,y_{i,<t})}
\]
改成 sequence-level ratio：
\[
s_i
=
\left[
\frac{\pi_\theta(y_i\mid x)}
{\pi_{\theta_\text{old}}(y_i\mid x)}
\right]^{1/|y_i|}
\]
然后做 sequence-level clipping 和 optimization。\\
核心直觉：reward 是给整段回答的，policy correction 也应该按整段回答来做。
\end{summarybox}

\subsubsection{本轮检查题}

\begin{enumerate}[leftmargin=1.5em]
  \item GRPO 的 token-level ratio 为什么和 sequence-level reward 粒度不匹配？
  \item GSPO 的 $s_i(\theta)$ 为什么要取 $1/|y_i|$ 次方？
  \item GSPO 的 objective 和 GRPO 的 objective 最大区别是什么？
  \item 为什么 GSPO 对 MoE RL 训练稳定性特别有帮助？
  \item GSPO 相比 GRPO 的 tradeoff 是什么？
\end{enumerate}
